\documentclass[11pt, a4paper, spanish, openright, twoside]{book}
\usepackage[spanish, activeacute]{babel}
\usepackage[utf8]{inputenc}
\usepackage[top=2.5cm, bottom=2.5cm, outer=1.75cm, inner=1.75cm, heightrounded, marginparwidth=2.5cm, marginparsep=0.3cm]{geometry}	%márgenes empequeñecidos
%\usepackage[top=2.5cm, bottom=2.25cm, outer=2.75cm, inner=2.75cm, heightrounded, marginparwidth=2.5cm, marginparsep=0.3cm]{geometry}	%márgenes originalmente
\usepackage{dpg}
\usepackage{afvc}
%Figuras
\usepackage[vflt]{floatflt}		%Entorno float-figure
%\graphicspath{/Temas/Tema01/Imagenes}

%%%%%%		chapter's style		%%%%%%%%%%%%%%%%%%%

\renewcommand{\thepage}{\arabic{page}}% Arabic page numbers\fancyhead{}
\pagestyle{fancy}
\fancyfoot{}
\fancyhead[RO]{}	%encabezado de pares: nombre de la sección
\fancyfoot[LE,RO]{\thepage}	%abajo a izqda en pares, derecha en impares: numero de pagina
\fancyhead[LE]{\nouppercase{\leftmark}} %cuadro izquierdo de pagina par: parte y contador
\fancyfoot[CE]{Tecnología y Organización de Computadores} 
\fancyfoot[CO]{Doble Grado Informática-Matemáticas - Universidad Complutense}
\renewcommand{\footrulewidth}{0.4pt}
\renewcommand{\headrulewidth}{0.4pt}		% linea por debajo del encabezado
\renewcommand{\sectionmark}[1]{\markright{\textbf{\thesection. #1}}}	%negrita
\renewcommand{\labelitemi}{$\circ$} %Primer itemize con circunferencia vacia
\renewcommand{\labelitemii}{$\cdot$} %Segundo itemize con punto pequeño
\renewcommand*{\thesection}{\arabic{section}}	% Hace que no apareca el indice de capitulos y que comience en section, GRACIAS A RUBEN
\setlength{\leftmarginii}{0em} %Segundo itemize sin sangria
\setlength{\leftmarginiii}{1em} %Tercer itemize casi sin sangria
\renewcommand{\labelitemiii}{ }
\pagenumbering{roman}
\newcommand*{\PKT}{\hbox{P}\kern-2.5pt\lower3.5pt\hbox{\small{K}}\kern-2.8pt\hbox{T}\kern-2pt}	%PiKey Team en bonito


\begin{document} 
\title{\Huge{\textsc{Manual Jetpack}} \\
	\vspace{0.7cm}
	 \textsc{\Large{Tecnología y Organización de Computadores}} \\
	\includegraphics[scale=0.3]{ucm.pdf}}
\author{{\Large{PiKey Team-}} \PKT \ : \vspace{0.2cm} \\
	Jesús Aguirre Pemán \\
	 Enrique Ballesteros Horcajo \\
	 Mayra Alexandra Castrosqui Florián\\
	 Jaime Dan Porras Rhee \\
	 Ignacio Iker Prado Rujas}
\date{\Today}
\maketitle

\newpage
\mbox{}
\thispagestyle{empty}						% Hoja en blanco, sin numeros ni nada
\newpage


\tableofcontents 							%INDICE hipervinculado

\newpage
\mbox{}
\thispagestyle{empty}						% Hoja en blanco, sin numeros ni nada
\newpage

\pagenumbering{arabic}						% Pone el contador de paginas a 1 y ahora en numeros normales

\vspace{3cm}

%ESTE FICHERO DEBERÍA SER DE UNAS DOS CARAS APROX.______________________________________________________

\section{Ficheros y dependencias}
El módulo central de nuestro proyecto es \textit{choca.vhd}. Es ahí dónde se gestionan las posiciones de Barry, el fondo, obstáculos y monedas. Sin embargo, para leer los datos de memoria utilizamos varias ROMs:
\begin{itemize}
\item \textit{rom_rgb_9b_fondo}: Memoria ROM que contiene los datos del fondo del juego.
\item \textit{rom_rgb_9b_game-over-negro}: Memoria ROM que contiene los datos del "Game Over".
\item \textit{rom_rgb_9b_Joyride}: Memoria ROM que contiene los datos del fondo del juego.
\item \textit{rom_rgb_9b_nivel_1_0}: Memoria ROM que contiene los datos del primer nivel del juego.
\item \textit{rom_rgb_9b_nivel_2_0}: Memoria ROM que contiene los datos del segundo nivel del juego.
\item OJO, PROBABLEMENTE HAYA ALGUNA MAS
\end{itemize}

Además, utilizamos distintas frecuencias para Barry, el fondo y los obstáculos, pues cada uno de ellos se mueve a una velocidad diferente. Para ello tenemos varios ficheros divisores:
\begin{itemize}
\item \textit{divisor.vhd}: Fichero divisor que controla la velocidad de ???.
\item \textit{divisor_choques}: Fichero divisor que controla la velocidad de ???.
\item \textit{divisor_movimiento_fondo}: Fichero divisor que controla la velocidad de movimiento del fondo.
\item \textit{divisor_movimiento_obstaculos}: Fichero divisor que controla la velocidad de movimiento de los obstáculos.
\item \textit{divisor_munyeco}: Fichero divisor que controla la velocidad de Barry.
\item OJO, PROBABLEMENTE HAYA ALGUNA MAS
\end{itemize}

Con el fichero \textit{contador.vhd} controlaremos el número de monedas cogidas, que se mostrarán en los displays 7 segmentos de la placa mediante el fichero \textit{pines.ucf}.

Mediante el ficher \textit{control_teclado.vhd} controlamos las teclas que pulsamos: la barra espaciadora para hacer ascender a Barry, y la 'p' para pausar la partida.

El ranndomgenerator no se si se usa ¿¿??


\section{nombre.bit}


\section{Pines y significado}
Conectaremos el teclado con los pines:

NET PS2CLK LOC = B16;
NET PS2DATA LOC = E13;
NET PS2CLK CLOCK_DEDICATED_ROUTE = FALSE;

Mostraremos el número de monedas mediante los displays de la placa extendida:
\begin{itemize}
\item placa izquierda:

NET displayizq<0> LOC=H14;
NET displayizq<1> LOC=M4;
NET displayizq<2> LOC=P1;
NET displayizq<3> LOC=N3;
NET displayizq<4> LOC=M15;
NET displayizq<5> LOC=H13;
NET displayizq<6> LOC=G16;

\item placa derecha:

NET displaydcha<0> loc=E2;
NET displaydcha<1> loc=E1;
NET displaydcha<2> loc=F3;
NET displaydcha<3> loc=F2;
NET displaydcha<4> loc=G4;
NET displaydcha<5> loc=G3;
NET displaydcha<6> loc=G1;
\end{itemize}

Por último, conectaremos los colores y la posición en la pantalla:

NET rgb<0> LOC=C9; #BLU0;
NET rgb<1> LOC=E7; #BLU1;
NET rgb<2> LOC=D5; #BLU2;
NET rgb<3> LOC=A8; #GRN0; 
NET rgb<4> LOC=A5; #GRN1;
NET rgb<5> LOC=C3; #GRN2;
NET rgb<6> LOC=C8; #RED0;
NET rgb<7> LOC=D6; #RED1;
NET rgb<8> LOC=B1; #RED2;
NET hsyncb LOC=B7; #HSYNC# 
NET vsyncb LOC=D8; #VSYNC# 

\section{Otros aquelarres de inicialización}
Qué carajos hay que poner aquí??


\section{Cómo se ``juega''}
Tras el despegue, simplemente pulsa la barra espaciadora para ascender y deja de pulsarla para descender. El objetivo es esquivar todos los obstáculos, y recoger el mayor número de monedas posible


\section{Qué espero ver}
La versión para FPGA del juego Jetpack Joyride. 


\end{document}